\documentclass[11pt,a4paper]{article}
% Shared preamble for RuPhO English problem statements (match T1 2026 style).
% Use: \input{latex/rupho_en_preamble.tex} after \documentclass.
\usepackage[margin=1in]{geometry}
\usepackage{amsmath,amssymb}
\usepackage{graphicx}
\usepackage{float}
\usepackage{enumitem}
\usepackage{microtype}
\usepackage{caption}
\usepackage{tabularx}
\usepackage{hyperref}

\captionsetup{font=small,labelfont=bf,skip=6pt}

\setlist[itemize]{leftmargin=1.2cm}
\setlist[enumerate]{leftmargin=1.2cm}

% One outer box per subproblem: [id | statement | points] (no inner rules)
\newcommand{\problembox}[3]{%
  \par\addvspace{0.42em}%
  \noindent
  \setlength{\fboxsep}{7.5pt}%
  \setlength{\fboxrule}{0.95pt}%
  \fbox{%
    \begin{minipage}{\dimexpr\linewidth-2\fboxsep-2\fboxrule\relax}
    \begin{tabularx}{\linewidth}{@{}>{\raggedright\arraybackslash\bfseries}p{2.1em} @{\hspace{0.9em}} >{\raggedright\arraybackslash}X @{\hspace{0.9em}} >{\raggedleft\arraybackslash\bfseries}p{2.4em}@{}}
    #1 & #2 & #3
    \end{tabularx}
    \end{minipage}%
  }%
  \par\addvspace{0.32em}%
}

\title{\textbf{Optics of water}}
\author{\normalsize Y24 (2024) --- English translation}
\date{}
\begin{document}
\maketitle

\section*{Introduction}

If you have gone swimming in a river at least once, you have most likely noticed the interesting optical properties of “real” water. We are talking mainly about the processes of scattering and absorption on the water itself and the structures inside it (particles, bubbles, etc.) [see. rice. 1].

\begin{figure}[H]
    \centering
    \includegraphics[width=0.72\textwidth]{PE_EN_assets/fig1.jpg}
    \caption{Figure 1.}
\end{figure}

These phenomena can tell an experienced researcher a lot about the properties and characteristics of water in a particular body of water, therefore methods for solving problems of scattering and absorption are also of applied interest. This is what this task will focus on.

\section*{Part A. Water molecule and absorption coefficient (3.5 points)}

The water molecule ($\mathrm{H_2O}$) consists of two hydrogen atoms and one oxygen atom, forming an isosceles triangle. Like any mechanical system, a water molecule can experience vibrations around its equilibrium position. All three possible modes of vibration of a water molecule are shown in Fig. 2.

\begin{figure}[H]
    \centering
    \includegraphics[width=0.52\textwidth]{PE_EN_assets/fig2.jpg}
    \caption{Figure 2.}
\end{figure}

When the frequency of an electromagnetic wave propagating in water turns out to be approximately a multiple of one of the natural frequencies of vibrations of the molecule, these vibrations are excited in the water molecules with a very large amplitude, “taking away” the energy of the wave and subsequently re-emitting it in all directions. This leads to the fact that the absorption coefficient has local maxima at these frequencies [see. rice. 3]. Note: The absorption coefficient determines the distance after which the beam intensity decreases by a factor of $e$. Dimension $[\kappa]=m^{-1}$.

\begin{figure}[H]
    \centering
    \includegraphics[width=0.72\textwidth]{PE_EN_assets/fig3.jpg}
    \caption{Figure 3.}
\end{figure}

Since the mass of the oxygen atom is much greater than the mass of the hydrogen atom, it can be considered practically motionless. Moreover, hydrogen atoms interact with each other noticeably less than with the oxygen atom. According to the results of experimental studies, the frequencies of the two vibration modes practically coincide, and the frequency of the third mode turns out to be two times less than the other two with good accuracy.

\problembox{A1}{Indicate which mode has the lowest natural frequency. Justify your answer.}{0.30}

\problembox{A2}{Find the wavelengths $\lambda$ and frequencies $\nu$ corresponding to the local maxima of the absorption coefficient.}{0.90}

It is clear that all frequencies corresponding to local maxima must be multiples of some fundamental frequency $\nu_0$. To find the fundamental frequency, it is necessary to “lay” these points on a straight line, i.e. select the frequency multiplicities so as to obtain a linear dependence. To do this, it is necessary to develop a criterion for “getting on the line.” Usually, hitting a straight line is taken as a criterion, taking into account the error cross. As an error, we take the half-width at half-height - the minimum frequency deviation at which the absorption indicator drops by 2 times.

\problembox{A3}{Find the half-width $\Delta\nu$ of the local maxima of the absorption coefficient.}{1.20}

\problembox{A4}{Based on the results obtained, determine the frequency multiplicity corresponding to the local maxima, and also find the main resonant frequency $\nu_0$ of vibrations of the water molecule. Graphically determine the error $\nu_0$.}{1.20}

\section*{Part B. Scattering by bubbles (6.5 points)}

In “real” water, not only absorption, but also scattering of light occurs. In this problem we consider scattering by air bubbles. An example is shown in Fig. 4.

\begin{figure}[H]
    \centering
    \includegraphics[width=0.72\textwidth]{PE_EN_assets/fig4.jpg}
    \caption{Figure 4.}
\end{figure}

For simplicity, we consider scattering by a single bubble [see rice. 5]. Let us introduce the relative impact parameter $\rho$ as the ratio of the impact parameter to the radius of the bubble, and denote the deflection angle as $\psi$. The refractive index inside the bubble is $1$, outside – $n$.

\begin{figure}[H]
    \centering
    \includegraphics[width=0.72\textwidth]{PE_EN_assets/fig5.jpg}
    \caption{Figure 5.}
\end{figure}

Then the deflection angle will be written as\textbackslash{}[\textbackslash{}psi=2\textbackslash{}left[\textbackslash{}arcsin n\textbackslash{}rho-\textbackslash{}arcsin\textbackslash{}rho\textbackslash{}right],\textbackslash{}]and the intensity of the scattered radiation - as\textbackslash{}[\textbackslash{}mathcal J(\textbackslash{}psi)\textbackslash{}equiv\textbackslash{}frac{\textbackslash{}rho\textasciitilde{} \textbackslash{}mathrm d\textbackslash{}rho}{\textbackslash{}sin\textbackslash{}psi\textasciitilde{}\textbackslash{}mathrm d\textbackslash{}psi}.\textbackslash{}]In Fig. Figure 6 shows a graph of the dependence $\mathcal J(\psi)$, obtained experimentally.

\begin{figure}[H]
    \centering
    \includegraphics[width=0.72\textwidth]{PE_EN_assets/fig6.jpg}
    \caption{Figure 6.}
\end{figure}

An honest solution to the inverse problem (determining the refractive index of a liquid $n$ from the distribution $\mathcal J(\psi)$) is possible, but quite problematic. As a compromise, it is proposed to consider the asymptotic behavior of the curve $\mathcal J(\psi)$ in a “convenient” approximation.

\problembox{B1}{Obtain an approximate expression for $\mathcal J(\psi)$, considering the angles $\psi$ small.}{1.50}

The table in the answer sheet shows the experimental points for $\mathcal J(\psi)$. $\mathcal J$ is given in conventional units.

\begin{table}[h]
\centering
\caption{Data: $\psi$ (degrees) vs $J$}
\label{tab:psi_j_data}
\begin{tabular}{cc|cc|cc|cc|cc}
\hline
$\psi$ (°) & $J$ & $\psi$ (°) & $J$ & $\psi$ (°) & $J$ & $\psi$ (°) & $J$ & $\psi$ (°) & $J$ \\
\hline
0.5 & 0.465 & 10.5 & 0.414 & 20.5 & 0.343 & 30.5 & 0.244 & 40.5 & 0.165 \\
1.0 & 0.464 & 11.0 & 0.409 & 21.0 & 0.343 & 31.0 & 0.250 & 41.0 & 0.166 \\
1.5 & 0.458 & 11.5 & 0.412 & 21.5 & 0.332 & 31.5 & 0.236 & 41.5 & 0.166 \\
2.0 & 0.462 & 12.0 & 0.417 & 22.0 & 0.321 & 32.0 & 0.231 & 42.0 & 0.171 \\
2.5 & 0.449 & 12.5 & 0.410 & 22.5 & 0.333 & 32.5 & 0.225 & 42.5 & 0.156 \\
3.0 & 0.444 & 13.0 & 0.409 & 23.0 & 0.317 & 33.0 & 0.235 & 43.0 & 0.165 \\
3.5 & 0.451 & 13.5 & 0.404 & 23.5 & 0.324 & 33.5 & 0.225 & 43.5 & 0.153 \\
4.0 & 0.459 & 14.0 & 0.393 & 24.0 & 0.315 & 34.0 & 0.212 & 44.0 & 0.154 \\
4.5 & 0.460 & 14.5 & 0.393 & 24.5 & 0.305 & 34.5 & 0.214 & 44.5 & 0.139 \\
5.0 & 0.445 & 15.0 & 0.391 & 25.0 & 0.299 & 35.0 & 0.221 & 45.0 & 0.148 \\
5.5 & 0.455 & 15.5 & 0.384 & 25.5 & 0.303 & 35.5 & 0.200 & 45.5 & 0.138 \\
6.0 & 0.444 & 16.0 & 0.391 & 26.0 & 0.300 & 36.0 & 0.197 & 46.0 & 0.145 \\
6.5 & 0.438 & 16.5 & 0.382 & 26.5 & 0.286 & 36.5 & 0.202 & 46.5 & 0.145 \\
7.0 & 0.448 & 17.0 & 0.378 & 27.0 & 0.282 & 37.0 & 0.194 & 47.0 & 0.143 \\
7.5 & 0.436 & 17.5 & 0.370 & 27.5 & 0.285 & 37.5 & 0.187 & 47.5 & 0.130 \\
8.0 & 0.434 & 18.0 & 0.362 & 28.0 & 0.275 & 38.0 & 0.188 & 48.0 & 0.132 \\
8.5 & 0.432 & 18.5 & 0.360 & 28.5 & 0.266 & 38.5 & 0.185 & 48.5 & 0.126 \\
9.0 & 0.442 & 19.0 & 0.352 & 29.0 & 0.264 & 39.0 & 0.180 & 49.0 & 0.119 \\
9.5 & 0.430 & 19.5 & 0.347 & 29.5 & 0.255 & 39.5 & 0.192 & 49.5 & 0.125 \\
10.0 & 0.429 & 20.0 & 0.340 & 30.0 & 0.259 & 40.0 & 0.181 & — & — \\
\hline
\end{tabular}
\end{table}

\problembox{B2}{Linearize the resulting dependence and find the refractive index of the liquid $n$. Note: It is not necessary to use all points to get a good result.}{5.00}

\vspace{1em}
\noindent\textit{Source:} \url{https://pho.rs/p/3864}

\end{document}